%%% FIELD five-node-quartic-statement
Let
\[
D=\{0\to1,0\to2,1\to3,2\to3,1\to4,2\to4\},\qquad
T=\{0\to1,0\to2,1\to3,1\to4\},
\]
let \(B=\{0\leftrightarrow i:1\le i\le4\}\), and query \(e=2\to4\). Then \((G,T,e)\in\mathcal G_{5,2}\) and
\[
d_e(G)=4.
\]
At the rational covariance
\[
\Sigma_\star=
\begin{pmatrix}
30&61&92&707&1014\\
61&155&187&1561&2247\\
92&187&314&2327&3332\\
707&1561&2327&17981&25747\\
1014&2247&3332&25747&36972
\end{pmatrix},
\]
the queried-coordinate fiber consists of four distinct real simple values, and every value extends to real \(\Lambda\) for which \(\Omega=(I_5-\Lambda)^{\mathsf T}\Sigma_\star(I_5-\Lambda)\) is positive definite and has the zero pattern required by \(B\). Appending any number of new children of vertex \(0\), while allowing every bidirected pair incident to each new child, preserves this original-coordinate degree. Consequently \(D_k(5)\ge4\) for every admissible \(k\ge2\).
%%% FIELD five-node-quartic-proof
Write
\[
a=\lambda_{01},\quad b=\lambda_{02},\quad c=\lambda_{13},\quad
d=\lambda_{23},\quad u=\lambda_{14},\quad f=\lambda_{24}.
\]
The set \(T\) has root \(0\), gives every other vertex exactly one parent, and has four edges. Hence it is a rooted spanning arborescence. Moreover \(D\setminus T=\{2\to3,2\to4\}\), so \((G,T,e)\in\mathcal G_{5,2}\) by \(\mathrm{def:graph\mbox{-}class}\).

First construct a regular real specialization. Set
\[
(a,b,c,d,u,f)=(2,3,4,5,6,7)
\]
and let \(\Omega_\star\) have diagonal \((30,31,32,33,34)\), entries \((\Omega_\star)_{0i}=i\) for \(1\le i\le4\), and all other off-diagonal entries zero. For every nonzero \(x\in\mathbb R^5\), the inequality \(2|x_ix_j|\le x_i^2+x_j^2\) gives
\[
x^{\mathsf T}\Omega_\star x
\ge \sum_{i=0}^4\left((\Omega_\star)_{ii}
-\sum_{j\ne i}|(\Omega_\star)_{ij}|\right)x_i^2>0,
\]
because the five parenthesized margins are \(20,30,30,30,30\). Thus \(\Omega_\star\succ0\). Direct multiplication in
\(\Sigma=(I_5-\Lambda)^{-\mathsf T}\Omega_\star(I_5-\Lambda)^{-1}\) gives exactly the displayed matrix \(\Sigma_\star\).

The missing bidirected pairs are precisely the six pairs inside \(\{1,2,3,4\}\). Therefore the specialization of the incidence ideal is generated in \(\mathbb Q[a,b,c,d,u,f]\) by
\[
\begin{aligned}
R_1={}&30ab-92a-61b+187,\\
R_2={}&61ac+92ad-707a-155c-187d+1561,\\
R_3={}&61au+92af-1014a-155u-187f+2247,\\
R_4={}&61bc+92bd-707b-187c-314d+2327,\\
R_5={}&61bu+92bf-1014b-187u-314f+3332,\\
R_6={}&155cu+187cf-2247c+187du+314df-3332d\\
&\hspace{2.8em}-1561u-2327f+25747.
\end{aligned}
\]
These are the entries of \((I_5-\Lambda)^{\mathsf T}\Sigma_\star(I_5-\Lambda)\) at those six missing pairs, so their vanishing is equivalent to the required off-diagonal zero pattern.

We next give the characteristic-zero elimination certificate. In lexicographic order
\(a>b>c>d>u>f\), exact Buchberger reduction over \(\mathbb Q\) gives the following basis; harmless nonzero rational normalizing constants have been suppressed:
\[
\begin{aligned}
H_1={}&384915840a+425450175590061f^3-8965411803407607f^2\\
&+62974399432228195f-147445028655847225,\\
H_2={}&741827520b-593375419233f^3+12321625507347f^2\\
&-85271405126447f+196665729339485,\\
H_3={}&615040c+25798931271f^3-543666188277f^2\\
&+3818873142057f-8941504654939,\\
H_4={}&4d-3f+1,\\
H_5={}&153760u+8599643757f^3-181222062759f^2\\
&+1272957714019f-2980501654153,\\
H_6={}&q(f)=(f-7)c(f),\\
c(f)={}&8599643757f^3-181150021425f^2
+1271943861499f-2976933819575.
\end{aligned}
\]
For completeness, this is a two-sided ideal certificate rather than a one-way eliminant: each \(H_j\) is obtained from the \(R_i\) by exact rational \(S\)-polynomial reductions, while direct division of every \(R_i\) by \((H_1,\ldots,H_6)\) has zero remainder. Hence
\[
(R_1,\ldots,R_6)=(H_1,\ldots,H_6).
\]
The leading monomials of the normalized \(H_j\) are respectively
\(a,b,c,d,u,f^4\); they are pairwise relatively prime, so Buchberger's product criterion also verifies that all remaining \(S\)-polynomials reduce to zero. This proves both ideal containments and the asserted lexicographic basis over characteristic zero.

The displayed basis identifies the specialized quotient with \(\mathbb Q[f]/(q)\), since \(H_1,\ldots,H_5\) recover exactly one value of \((a,b,c,d,u)\) from each value of \(f\). Exact integer arithmetic gives
\[
c(7)=-30256
\]
and
\[
\operatorname{disc}(c)
=1802165059445823134111101206528>0.
\]
The nonzero discriminant says that the cubic has no repeated root. A real cubic has either three real roots or one real root and a nonreal conjugate pair; from the product formula for its discriminant, the latter case has negative discriminant. Thus \(c\) has three distinct real roots. Since \(c(7)\ne0\), these roots and \(7\) are four distinct real roots of \(q\). It follows that the specialized incidence fiber is a reduced set of four points with pairwise distinct original-coordinate values \(f=\lambda_{24}\).

We now pass from the specialization to the generic coordinate degree. Partition the six coefficient variables into the four incoming-vertex blocks
\[
\{a\},\qquad\{b\},\qquad\{c,d\},\qquad\{u,f\}.
\]
The six residuals have binary block graph \(K_4\) and degree one in each incident block. The multihomogeneous coefficient in \(\mathrm{lem:multihomogeneous\mbox{-}bezout\mbox{-}isolated\mbox{-}bound}\) is
\[
[X_1X_2X_3^2X_4^2]\prod_{1\le i<j\le4}(X_i+X_j)=4.
\]
Indeed, in the factor order \(12,13,14,23,24,34\), the only four choices of a variable with exponent vector \((1,1,2,2)\) are
\[
\begin{gathered}
(1,3,4,2,4,3),\qquad(1,3,4,3,2,4),\\
(2,1,4,3,4,3),\qquad(2,3,1,3,4,4).
\end{gathered}
\]
Consequently every isolated generic solution, counted with multiplicity, contributes to a total of at most four.

It remains to justify that the four specialized points are genuine generic branches. Since their fiber algebra is reduced and zero-dimensional, its cotangent space at every point is zero. The six \(R_i\) generate the fiber ideal in six variables, so their Jacobian with respect to \((a,b,c,d,u,f)\) has rank six at each point. Equivalently, this also follows by differentiating the triangular basis above, whose diagonal derivatives in \((a,b,c,d,u)\) are nonzero and whose last derivative is \(q'(f)\ne0\). For covariance deviations \(z\), expand the residual equations degree by degree in a putative series \(\lambda(z)\). At degree one the invertible Jacobian determines the coefficient uniquely; after degrees below \(r\) have been determined, the degree-\(r\) equation has the same Jacobian and hence uniquely determines the degree-\(r\) coefficient. This constructs four distinct formal solution branches over \(\mathbb C[[z]]\), and their \(f\)-series remain distinct because their constant terms are distinct.

The graph has six directed coefficients, five disturbance variances, and four supported off-diagonal disturbance coordinates, hence fifteen source coordinates, equal to the fifteen covariance coordinates. The existence of any one full-rank formal inverse makes the covariance parametrization dominant; thus its covariance function field is embedded by the independent covariance deviations. By \(\mathrm{thm:incidence\mbox{-}localization}\), the generic queried-coordinate image is the root set of the minimal polynomial of \(f\). The four formal branches give four distinct roots of that polynomial, so \(d_e(G)\ge4\). The multihomogeneous bound just proved gives \(d_e(G)\le4\), and therefore
\[
d_e(G)=4.
\]
This argument concerns the original coefficient \(\lambda_{24}\), not merely a primitive element for the whole fiber.

Every one of the four roots is real, and the five triangular recovery equations have real coefficients and nonzero constant multipliers; hence every root gives a real matrix \(\Lambda_j\). Define
\[
\Omega_j=(I_5-\Lambda_j)^{\mathsf T}\Sigma_\star(I_5-\Lambda_j).
\]
Acyclicity makes \(I_5-\Lambda_j\) unit triangular in the order \(0,1,2,3,4\), hence invertible. For nonzero real \(x\),
\[
x^{\mathsf T}\Omega_jx
=\bigl((I_5-\Lambda_j)x\bigr)^{\mathsf T}
\Sigma_\star\bigl((I_5-\Lambda_j)x\bigr)>0,
\]
because \(\Sigma_\star\succ0\), itself following from its congruence with \(\Omega_\star\succ0\). Thus \(\Omega_j\succ0\). The six residual equations say exactly that all off-diagonal entries missing from \(B\) vanish, so all four real roots have supported positive-definite lifts.

Finally append one new vertex \(w\) as a child of \(0\), put \(0\to w\) into \(T\), and allow every bidirected pair incident to \(w\). There is then no new missing-pair residual involving \(w\). Over the old source field, congruence by \(I-\Lambda\) is an invertible linear change between the new disturbance row and the new covariance row; the new arrow coefficient is an additional free transcendental. Hence the enlarged source field can be written
\[
L'=L(\lambda_{0w},\sigma_{iw}:i\in V,\sigma_{ww}),
\]
while the enlarged covariance field is
\[
K'_G=K_G(\sigma_{iw}:i\in V,\sigma_{ww}),
\]
with all displayed adjoining variables algebraically independent over the respective old fields. Irreducibility of the old minimal polynomial is preserved by a purely transcendental base extension, so
\[
[K'_G(\lambda_e):K'_G]=[K_G(\lambda_e):K_G]=4.
\]
Iteration proves the appending claim. For \(n=5\), the same graph belongs to \(\mathcal G_{5,k}\) whenever admissible \(k\ge2\), since its excess is two. The definition of \(D_k(5)\) therefore yields \(D_k(5)\ge4\).
%%% FIELD five-node-census-proposed-change
Using \(\mathsf{CensusHandle}\), complete the secondary exhaustive characteristic-zero coordinate census for every rooted graph-query orbit \((G,T,e)\in\mathcal G_{n,k}\) with \(n\le5\) and \(k\le2\). For each orbit compute \(d_e(G)\) from \(A_K=S_G^{-1}(R_{\mathrm{inc}}/I_G)\), provide a saturated rational-univariate certificate for every finite coordinate fiber, and compare the resulting degree with its prescribed-orientation coefficient, explicitly recording every strict loss caused by forced covariance relations, affine dehomogenization, or projective-infinity roots. Treat the proved five-node cubic coordinate case and the proved five-node quartic with four real positive-definite lifts as fixed base cases. The census is a regression and structural-classification problem; it is not a prerequisite for the already determined numerical frontier \(D_k(n)\) when \(n\le5\).
%%% FIELD five-node-census-statement-replace
Using \(\mathsf{CensusHandle}\), complete the secondary exhaustive characteristic-zero coordinate census for every rooted graph-query orbit \((G,T,e)\in\mathcal G_{n,k}\) with \(n\le5\) and \(k\le2\). For each orbit compute \(d_e(G)\) from \(A_K=S_G^{-1}(R_{\mathrm{inc}}/I_G)\), provide a saturated rational-univariate certificate for every finite coordinate fiber, and compare the resulting degree with its prescribed-orientation coefficient, explicitly recording every strict loss caused by forced covariance relations, affine dehomogenization, or projective-infinity roots. Treat the proved five-node cubic coordinate case and the proved five-node quartic with four real positive-definite lifts as fixed base cases. The census is a regression and structural-classification problem; it is not a prerequisite for the already determined numerical frontier \(D_k(n)\) when \(n\le5\).
%%% FIELD five-node-census-partial
The exact rational certificate above proves that the specified graph in \(\mathcal G_{5,2}\) has \(d_{2\to4}(G)=4\), with four simple real queried values and four supported positive-definite lifts. Together with the canonical cubic transfer and small-size frontier results, the fixed numerical base cases are
\[
D_0(2)=1,\qquad D_k(3)=1,\qquad D_k(4)=2,
\]
for every admissible \(k\), and
\[
D_0(5)=2,\qquad D_1(5)=3,\qquad D_k(5)=4\quad(2\le k\le6).
\]
What remains is the finite per-orbit coordinate classification and its saturated certificates, not determination of these maxima.
